\documentclass[10pt,slidestop,usepdftitle=false]{beamer}
\usepackage[accumulated]{beamerseminar}
\usepackage{beamertexpower}
\usepackage{beamerthemeshadow}
\usepackage[T2A]{fontenc}       %поддержка кириллицы
\usepackage[utf8]{inputenc}   %пока бибтех не дружит до конца с юникодом
\usepackage[russian]{babel}     %определение языков в документе
\usepackage{amssymb,amsmath}    %математика
\usepackage{mathtools}
\usepackage{multirow}
\usepackage{amsmath}
\usepackage{multicol}
\usepackage{graphicx}
\usepackage{float}
\usepackage{subcaption}
\newcommand{\eex}{\end{example}}
\newcommand{\bea}{\begin{eqnarray*}}
\newcommand{\eea}{\end{eqnarray*} }
\newcommand{\be}{\begin{eqnarray}}
\newcommand{\ee}{\end{eqnarray} } 
\DeclareMathOperator*{\argsup}{arg\,sup}
\DeclareMathOperator*{\arginf}{arg\,inf}
\DeclareMathOperator*{\argmax}{arg\,max}
\DeclareMathOperator*{\argmin}{arg\,min}

\title[Линейный и нелинейный регрессионный анализ] 
{Линейный и нелинейный регрессионный анализ}
\author{}
\date{}

\begin{document}

\newtheorem{defin}{Определение}
\newtheorem{mytheorem}{Теорема}
\newtheorem{mylemma}{Лемма}
\newtheorem{proposition}{Утверждение}
\newtheorem{algorithm}{Алгоритм}
\newtheorem{rem}{Замечание}

\begin{frame}
  \maketitle
\end{frame}

\begin{frame}{Содержание}
 \begin{itemize}
 
 \bigskip 

 \item

{\bf Введение}

\bigskip 
\bigskip
\pause    
 \item
 
{\bf 
 Постановка задачи регрессионного анализа} 
\bigskip 
\bigskip 
\pause 
 \item
 
 
{\bf Линейные модели. Свойства оценок по методу наименьших квадратов} 
 
 \bigskip 
\bigskip 
\pause
 \item
 
 
 {\bf Доказательство
теоремы Гаусса-Маркова}

\bigskip 
\bigskip 
\pause 
 \item

{\bf Нелинейные модели. Асимптотические свойства оценок параметров} 


\end{itemize}
\end{frame}
\begin{frame}
Началом систематического развития теории планирования эксперимента явились работы выдающегося английского статистика сэра Рональда Фишера, вошедшие в его широко известную монографию (Fisher, 1935). 

\bigskip
\pause

В этих работах комбинаторные методы применялись для выбора уровней факторов с целью повышения информативности эксперимента в задачах дисперсионного анализа. Первоначальным мотивом этих исследований была проблема повышения урожайности сельскохозяйственных культур.

\bigskip
\bigskip
\pause

 Однако глубокие математические результаты, полученные Фишером, послужили источником для развития самостоятельной ветви математической статистики, имеющей свою внутреннюю логику и широчайший спектр приложений.

\end{frame}
\begin{frame}
В 50-е годы XX века, в работах американских математиков Г. Элвинга (Elfving, 1952) и Г. Чернова (Chernoff, 1953) возникло новое направление, получившее название теории оптимального планирования эксперимента. 
\bigskip

\pause

С конца 50-х годов эта теория систематически развивалась в работах американских математиков Дж. Кифера и Дж. Вольфовица, а с конца 60-х значительный вклад к ее развитие внесли российские математики В.В. Федоров, Г.К. Круг и многие другие.

\bigskip

\pause

В основе этой теории лежит понятие плана эксперимента как дискретной вероятностной меры, а под оптимальным планом понимается план, доставляющий экстремальное значение некоторой выпуклой (или вогнутой) функции, заданной на множестве информационных матриц.

\bigskip
\pause

Родственное направление, связанное с исследованием поверхности отклика для нахождения его экстремальных значений, было инициировано в работах американского статистика Дж. Бокса  В России это направление развивалось в работах В.В. Налимова и его учеников (первая российская книга на эту тему – Налимов, Чернова, 1965).


\end{frame}
\begin{frame}{Постановка задачи регрессионного анализа}
 Рассмотрим уравнение:
 \begin{align*}
	y_{i} = \eta(x_i,\theta) + \varepsilon_{i}, \qquad &i = 1, \dots, N,
 \end{align*}
 где 
 \begin{itemize}
 \item $y_1,\dots,y_N \in \mathbb{R}$ --- результаты наблюдений,
 \pause
 \item $x_1,\dots,x_N \in \mathcal{X}$ --- условия проведения эксперимента (могут повторяться),
 \pause
 \item $\mathcal{X}$- некоторое заданное множество, например, ограниченное замкнутое подмножество конечномерного евклидова пространства,
 \pause
 \item Функция $\eta(x,\theta) : \mathcal{X} \times \Theta \rightarrow \mathbb{R}$ называется регрессионной моделью или функцией регрессии,
 \pause
 \item $\theta = (\theta_1,\dots,\theta_m)^T \in \Theta$ --- вектор неизвестных параметров,
 \pause
 \item $\varepsilon_1, \dots, \varepsilon_N$ --- случайные ошибки.
 \end{itemize}
 Это уравнение называется уравнением регрессии.
\end{frame}

\begin{frame}
Уже в простейшем случае,когда
 \begin{align*}
\eta(x,\theta)=\theta_1 +\theta_2 x
 \end{align*}
уравнение регрессии имеет многочисленные применения. Например, в экономике оно описывает соотношение спрос - предложение.
\pause
\bigskip 

Некоторые конкретные модели будут подробно рассмотрены в данном курсе. 
\pause
\bigskip

Основные понятия, относящиеся к линейным по параметрам моделям, можно найти в учебном пособии 

\bigskip

В.Б.Мелас, П.В.Шпелев "Планирование и анализ для регрессионных моделей", СПб, изд-во СПбГУ,2014.
\pause
\bigskip

По нелинейным моделям см. В.Б.Мелас "Локально оптимальные планы эксперимента". СПб, изд-во СПбГТУ, 1999.
\end{frame}

\begin{frame}
Уравнение регрессии:
 \begin{align*}
	y_{i} = \eta(x_i,\theta)+\varepsilon_{i}, \qquad &i = 1, \dots, N.
 \end{align*}
 \pause{Два разных типа задач по оценке параметров:}
 \begin{enumerate}
 \item<3-> Дана совместная выборка $(x_1, y_1), \dots, (x_N, y_N)$ и некоторая априорная информация о виде ошибок наблюдения $\varepsilon_1, \dots,\varepsilon_N$. 
 
 \vspace{2mm}
 
 Необходимо найти функцию $\eta(x,\theta)$, которая наилучшим с некоторой точки зрения образом аппроксимирует экспериментальные данные, и оценить ее параметры.

 \vspace{2mm} 
 
 \item<4-> Дана функция $\eta(x,\theta)$ с точностью до неизвестных параметров $\theta$ и некоторая априорная информация о виде ошибок наблюдения $\varepsilon_1, \dots,\varepsilon_N$. 
  
  \vspace{2mm}
 
 Необходимо таким образом выбрать условия проведения эксперимента $x_1,\dots,x_N$, чтобы по получаемой впоследствии выборке $(x_1, y_1), \dots, (x_N, y_N)$ можно было бы наилучшим с некоторой точки зрения образом оценить параметры  $\theta$.
 \end{enumerate}
\end{frame}

\begin{frame}
Уравнение регрессии:
 \begin{align*}
	y_{i} = \eta(x_i,\theta) + \varepsilon_{i}, \qquad &i = 1, \dots, N.
 \end{align*}
 Под планом эксперимента можно понимать совокупность не обязательно различных точек $x_1, \dots, x_N$, представлящих собой условия проведения экспериментов.  

\pause
Однако удобнее рассматривать план как дискретную вероятностную меру, приписывая каждой точке "вес" $1/N$.
\begin{defin}
Точный план эксперимента --- дискретная вероятностная мера 
\begin{align*}
\xi_N =
\begin{bmatrix}
	x_1 & \dots & x_N \\ 
	{1}/{N} & \dots & {1}/{N}
\end{bmatrix},
\end{align*}
где $x_1,\dots,x_N \in \mathcal{X}$ --- не обязательно различны. 
\end{defin}
\pause
Будем пока считать, что точный план эксперимента фиксирован. Разберемся с тем, как оценивать параметры модели при фиксированном плане. 
\end{frame}

\begin{frame}{Линейные модели. Свойства оценок по методу наименьших квадратов }
Уравнение регрессии:
 \begin{align*}
	y_{i} = \eta(x_i,\theta) + \varepsilon_{i}, \qquad &i = 1, \dots, N.
 \end{align*}
 Наши предположения
 \begin{enumerate}
 \item Случайные ошибки $\varepsilon_1,\dots,\varepsilon_N$ обладают свойствами:   
 
 $\mathbf{E}\varepsilon_i = 0$,
 $\mathbf{E}\varepsilon_i^2 = \sigma^2$, 
 $\mathbf{E}\varepsilon_i\varepsilon_j = 0,\ (i\ne j)$.
 
 \item  Параметры модели вещественные, т.е. $\Theta = \mathbb{R}^m$.
 \end{enumerate}
\pause 

Регрессионная модель называется линейной (по параметрам), если она может быть записана в виде
 \begin{align*}
\eta(x,\theta) = \theta_1 f_1(x) + \dots + \theta_m f_m(x),
\end{align*}где
 \begin{align*}
 f_1(x), \dots , f_m(x)-
\end{align*}
заданные функции, например, $f_i= x^{i-1},i=1,2,\dots,m$
(в этом случае модель называется полиномиальной).

\end{frame}
 \begin{frame}
Введем обозначения
\begin{align*}
Y =\left[\begin{array}{c}
                     y_1  \\
                     \cdots \\
                     y_N \\
                   \end{array}\right],\ \ \
                   X= \left[\begin{array}{ccc}
                     f_1(x_1) & \ldots & f_m(x_1) \\
                     \cdots & \cdots & \cdots \\
                     f_1(x_N) & \ldots & f_m(x_N) \\
                   \end{array}\right],  \ \ \ 
                   \varepsilon = \left[\begin{array}{c}
                     \varepsilon_1  \\
                     \cdots \\
                     \varepsilon_N \\
                   \end{array}\right]
\end{align*}
Для линейной модели уравнение регрессии переписывается как
\begin{align*}
Y = X \theta + \varepsilon.
\end{align*}
\end{frame}
\begin{frame}

\pause
Гауссом был предложен метод наименьших квадратов.

\pause
\bigskip
Оценка вектора параметров $\theta$ по методу наименьших квадратов определяется
следующим образом
\begin{align*}
\widehat{\theta} = \argmin_{\theta \in R^{m}}\left(Y-X\theta\right)^T\left(Y-X\theta\right).
\end{align*}

\pause
\bigskip
Справедлива следующая лемма.

\begin{lemma}{Лемма 1.}
Для любой матрицы $X$ и вектора $Y$ соответствующих
размерностей, система уравнений
\begin{eqnarray}
X^TX\theta=X^TY, \label{SPV:equat17}
\end{eqnarray}
имеет решение. Причем вектор $\theta^*$, удовлетворяющий системе,
является МНК-оценкой. 
\label{lem1.1.1} 
\end{lemma}
\end{frame}
\begin{frame}
{\bf Доказательство леммы \ref{lem1.1.1}}.

Для произвольной матрицы А векторное пространство, порожденное ее
столбцами, обозначим через $\mathcal{L}(A).$
\bigskip
\pause

 Покажем, что
\begin{eqnarray}
\mathcal{L}(A^T)=\mathcal{L}(A^TA). \label{SPV:equat18}
\end{eqnarray}
Действительно, для любого вектора $b,$ ортогонального пространству
$\mathcal{L}(A^T)$, имеет место равенство $b^TA^T=\textbf{0}
\Rightarrow b^TA^TA=\textbf{0}$.
\bigskip
\pause

В обратную сторону, если вектор $b$ ортогонален пространству
$\mathcal{L}(A^TA)$, т.е. $b^TA^TA=\textbf{0},$ то выполнено
$\textbf{0}=b^TA^TAb=(Ab)^TAb$ и, следовательно, $bA^T=\textbf{0}$.
\pause
\bigskip

Из равенства \eqref{SPV:equat18} автоматически следует, что для
любого вектора $X^TY \in \mathcal{L}(X^T)$ существует вектор
$X^TX\theta^* \in \mathcal{L}(X^TX)$ такой, что $X^TX\theta^*=X^TY$.
\end{frame}
\begin{frame}
Покажем теперь, что любое решение уравнения \eqref{SPV:equat17}
является МНК-оценкой, т.е. минимизирует функцию
\begin{eqnarray*}
\left(Y-X\theta\right)^T\left(Y-X\theta\right)
\end{eqnarray*}
\pause
\bigskip

Обозначим через $\widehat{\theta}$ произвольное решение
\eqref{SPV:equat17}, имеем:
\begin{eqnarray*}
\left(Y-X\theta\right)^T\left(Y-X\theta\right)=\\
\left(Y-X\widehat{\theta}+X(\widehat{\theta}-\theta)\right)^T\left(Y-X\widehat{\theta}+X(\widehat{\theta}-\theta)\right)\stackrel{*}{=}\\
\left(Y-X\widehat{\theta}\right)^T\left(Y-X\widehat{\theta}\right)+(\widehat{\theta}-\theta)^TX^TX(\widehat{\theta}-\theta)\geq\\
\left(Y-X\widehat{\theta}\right)^T\left(Y-X\widehat{\theta}\right).
\end{eqnarray*}

\pause

Равенство $"*"$ выполнено в силу того, что $\widehat{\theta}$
является решением уравнения \eqref{SPV:equat17} и, следовательно,
\begin{eqnarray*}
(\widehat{\theta}-\theta)^TX^T\left(Y-X\widehat{\theta}\right)=(\widehat{\theta}-\theta)^T\left(X^TY-X^TX\widehat{\theta}\right)=0
\end{eqnarray*}

\
 Таким образом, минимум
$\left(Y-X\theta\right)^T\left(Y-X\theta\right)$ достигается при
$\theta=\widehat{\theta}$. Лемма доказана. $$ \eqno \Box$$
\end{frame}
\begin{frame}
Заметим,что если матрица $X^TX$ невырожденная, то система
\eqref{SPV:equat17} имеет единственное решение:
\begin{eqnarray}\label{LSE}
\widehat{\theta}=(X^TX)^{-1}X^TY
\end{eqnarray}

\bigskip
\pause

Требования, предъявляемые к оценкам:
\begin{enumerate}
\bigskip
\pause

\item $\mathbf{E} \widehat\theta = \theta$ --- несмещенность,
\pause
\bigskip

\item $\widehat\theta = S Y$ --- линейность,
\bigskip
\pause

\item $\forall z \in \mathbb{R}^m$, $\forall \tilde{\theta} : \mathbf{E}\tilde{\theta} = \theta$, $\mathbf{D}z^T(\widehat\theta-\theta) \leq \mathbf{D}z^T(\tilde\theta-\theta)$ --- минимальность дисперсии.
\end{enumerate}
\bigskip
\pause

Будем предполагать, что матрица $X^T X$ - не вырожденная.

\pause
\begin{mytheorem}[Гаусса-Маркова]
Если матрица $X^T X$ - не вырожденная, то оценка по методу наименьших квадратов~(\ref{LSE}) является наилучшей (в смысле требования 3) оценкой в классе линейных несмещенных оценок.

Кроме того  , ковариационная матрица оценки имеет вид:
$$
\mathbf{D}\widehat\theta = \sigma^2 \left( X^T X \right)^{-1}.
$$
\end{mytheorem}
\end{frame}
\begin{frame}
Рассмотрим для примера линейную по параметрам модель
\begin{eqnarray*}
y_j = a +
bt_j+\epsilon_j,\ \ j=1,\ldots, N,
\end{eqnarray*}
где $t_1,\ldots,t_N$ --- некоторые заданные вещественные числа.
\begin{eqnarray*}
X= \left( \begin{array}{cc}
          1 & t_1 \\
          \vdots & \vdots \\
          1 & t_N \\
          \end{array}
          \right), \ \ X^TX= \left( \begin{array}{cc}
          N & \sum_{i}t_i \\
          \sum_{i}t_i & \sum_{i}t_i^2 \\
          \end{array}
          \right).
\end{eqnarray*}

\pause
Применяя известную формулу для обратной матрицы, получаем
\begin{eqnarray*}
(X^TX)^{-1}= \left( \begin{array}{cc}
          \frac{\sum_{i}t_i^2}{\Delta}  &\frac{ -\sum_{i}t_i}{\Delta} \\
         \frac{ -\sum_{i}t_i}{\Delta} & \frac{N}{\Delta} \\
          \end{array}
          \right),\ \ \Delta= N\sum_{i}t_i^2-\left(\sum_{i}t_i\right)^2.
\end{eqnarray*}

\pause
 
Искомая оценка
\begin{eqnarray*}
      \left(
       \begin{array}{c}
         \widehat{a} \\
         \widehat{b} \\
       \end{array}
              \right)=(X^TX)^{-1}X^TY=\left(
                              \begin{array}{c}
                                \frac{\sum_{i}t_i^2\sum_{i}y_i-\sum_{i}t_i\sum_{i}y_it_i}{\Delta} \\
                                \frac{N\sum_{i}y_it_i-\sum_{i}t_i\sum_{i}y_i}{\Delta} \\
                              \end{array}
                            \right).
\end{eqnarray*}

\end{frame}
\begin{frame}{Доказательство
теоремы Гаусса-Маркова.}

 Для доказательства нам
потребуются следующие вспомогательные утверждения:
\begin{lemma}{Лемма 2.}
МНК-оценка является линейно несмещенной оценкой (т.е. для нее
выполнены требования (a)-(b)).
 \label{lemMSE1}
\end{lemma}
{\bf Доказательство леммы \ref{lemMSE1}} Линейность МНК-оценки
очевидна ($S=(X^TX)^{-1}X^T$), проверим несмещенность. Посчитаем
математическое ожидание:
\begin{eqnarray*}
E\widehat{\theta}=ESY=SEY=SE(X\theta+\epsilon)=SX\theta=
(X^TX)^{-1}X^TX\theta=\theta
\end{eqnarray*}
 Лемма доказана. $$ \eqno \Box$$
 \end{frame}
\begin{frame}
Теперь сформулируем более общее утверждение о несмещенности оценки.
\begin{lemma}{Лемма 3.}
Для того, чтобы оценка $\widetilde{\theta}=AY$ была несмещенной,
необходимо и достаточно, чтобы выполнялось равенство:
\begin{eqnarray*}
AX=I
\end{eqnarray*}
 \label{lemMSE2}
\end{lemma}
Эта лемма является основным инструментом доказательства
теоремы Гаусса-Маркова.

\pause

{\bf Доказательство леммы \ref{lemMSE2}} $\newline$
\textbf{Достаточность.} Пусть $AX=I,$ тогда
\begin{eqnarray*}
E\widetilde{\theta}=EAY= \\
AE(X\theta+\epsilon)=\underbrace{AX}_{I}\theta=\theta.
\end{eqnarray*}
Таким образом, оценка $\widetilde{\theta}=AY$ является несмещенной оценкой.
\end{frame}
\begin{frame}
\textbf{Необходимость.} Пусть, теперь, выполнено
$E\widetilde{\theta}=\theta$, т.е. оценка несмещенная. Тогда получаем
$E\widetilde{\theta}= AX\theta=\theta$. Следовательно,
$$
AX\theta=\theta.
$$
\bigskip
\pause
Так
как это равенство должно быть выполнено для любого вектора
$\theta$, можем в качестве него выбрать вектор
$e_i=(0,\ldots,0,1,0,\ldots,0)^T$ с единицей на $i$-м месте.

\bigskip
\pause
Проделав это для всех $i=1,\ldots,m$, получим систему уравнений
$AX=I$, что завершает доказательство.
$$ \eqno \Box$$
\end{frame}
\begin{frame}
\begin{lemma}{Лемма 4.}
В предположениях теоремы Гаусса - Маркова дисперсионная матрица любой
линейной несмещенной оценки $\widetilde{\theta}=AY$ имеет вид:
\begin{eqnarray*}
D_{\widetilde{\theta}}=\sigma^2AA^T.
\end{eqnarray*}
В частном случае, если $\widetilde{\theta}$ - МНК-оценка, то
\begin{eqnarray*}
D_{\widetilde{\theta}}=\sigma^2(X^TX)^{-1}.
\end{eqnarray*}
 \label{lemMSE3}
\end{lemma}

\pause

{\bf Доказательство леммы \ref{lemMSE3}} вычислим дисперсионную матрицу оценки
$\widetilde{\theta}$:
\begin{eqnarray*}
D_{\widetilde{\theta}}=E(\widetilde{\theta}-\theta)(\widetilde{\theta}-\theta)^T=
E(AY-\theta)(AY-\theta)^T=EAYY^TA^T-\theta\theta^T,
\end{eqnarray*}
так как $EAY=(EY^TA^T)^T=\theta$. 
\pause
С другой стороны,
\begin{eqnarray*}
EAYY^TA^T=AE(YY^T)A^T=AE(X\theta+\epsilon)(X\theta+\epsilon)^TA^T=\\
\underbrace{AX}_I\theta\theta^T{\underbrace{X^TA}_I}^T+\sigma^2AA^T=\theta\theta^T+\sigma^2AA^T
\end{eqnarray*}

\pause

Таким образом, получаем, что $D_{\widetilde{\theta}}=\sigma^2AA^T$
и, в частности, в случае МНК-оценки (т.е. при $A=(X^TX^T)^{-1}X$)
\begin{eqnarray*}
D_{\widehat{\theta}}=\sigma^2AA^T=\sigma^2(X^TX)^{-1}X^TX(X^TX)^{-1}=\sigma^2(X^TX)^{-1}.
\end{eqnarray*}
Лемма доказана. $$ \eqno \Box$$ 
\end{frame}
\begin{frame}
Вернемся теперь к доказательству
теоремы. 
Проверим, что выполняется неравенство:
\begin{eqnarray*}
D_{\widehat{\theta}}\leq D_{\widetilde{\theta}}.
\end{eqnarray*}
\pause
Обозначим матрицу $(X^TX)^{-1}X^T$ через S. Для произвольной
несмещенной оценки $\widetilde{\theta}=AY$ имеем (в силу Леммы 4)
\begin{eqnarray*}
D_{\widetilde{\theta}}=\sigma^2 AA^T,
\end{eqnarray*}
а для оценки метода наименьших квадратов
\begin{eqnarray*}
D_{\widehat{\theta}}=\sigma^2 SS^T,
\end{eqnarray*}.
Значит нам достаточно доказать, что 
\begin{eqnarray*}
AA^T\geq  SS^T.
\end{eqnarray*}
\pause
Заметим, что для произвольной матрицы $L$
\begin{eqnarray*}
LL^T\geq 0,
\end{eqnarray*}
где 0 - нулевая матрица.
Подставляя сюда $L= A-S$ и произведя вычисления с учетом того,что $AX=I$,
получим 
$$
LL^T= AA^T-SS^T\geq 0,\, AA^T\geq SS^T.
$$
%откуда следует требуемое неравенство.
%Теорема  доказана. $$ \eqno \Box$$
\end{frame}
\begin{frame}
{\bf МНК-оценка в вырожденном случае.}

В предыдущем параграфе мы рассматривали задачу нахождения МНК-оценки
в случае, когда матрица $X^TX$ невырожденная. В данном параграфе мы
обобщим полученные результаты для 
когда ранг матрицы $X^TX$ меньше числа оцениваемых параметров. 
\bigskip
\pause

Заметим, что  матрица $X^TX$ является вырожденной, напрример,  в задачах дисперсионного анализа.
\bigskip
\pause

Итак, пусть $X^TX$ - вырожденная матрица. В этом случае, под
МНК-оценкой,
 по-прежнему, понимается вектор, минимизирующий $\|Y-X\theta \|^2$ и,
 согласно лемме \ref{lem1.1.1}, являющийся решением нормального уравнения
 \begin{eqnarray*}
X^TX\theta=X^TY.
\end{eqnarray*}
\bigskip
\pause

Как уже было показано, эта cистема уравнений совместна. Однако, в
вырожденном случае, она имеет бесконечное множество решений, т.е.
вектор $\widehat{\theta},$ являющийся её решением, определен не однозначно. 
\end{frame}
\begin{frame}
 
 

Пусть имеется произвольная система уравнений $Ax=y,$
где $x\in R^{n},y\in R^{m}$- вектора , $A\in R^{n\times m}$ -
матрица соответствующей размерности (не обязательно квадратная).
\bigskip
\pause

{\bf Определение. Обобщенно-обратной для матрицы $A$, называется матрица порядка
$n\times m,$ обозначаемая $A^{-}$, такая, что для любого вектора
$y$, для которого система $Ax=y$ совместна, вектор $x=A^{-}y$
является ее решением.} 
\pause 
\bigskip

 Для начала, отметим одно важное свойство,
которым обладает матрица $A^{-}$, а потом ответим на вопрос о ее
существовании.
\begin{lemma}{Лемма 5.}
Для того, чтобы матрица $B$ являлась обобщенно-обратной матрицей для
матрицы $A$, необходимо и достаточно, чтобы $B$ удовлетворяла
равенству:
\begin{eqnarray*}
ABA=A.
\end{eqnarray*}
 \label{lemGIM1}
\end{lemma}
\end{frame}
\begin{frame}

{\bf Доказательство леммы \ref{lemGIM1}}

Пусть для матрицы $A$ существует обобщенно-обратная матрица
$B=A^{-}$. Выберем в качестве $y$ $i$-ый столбец матрицы $A$.
Система
\begin{eqnarray*}
Ax=a_i,
\end{eqnarray*}
очевидно, совместна. Ее решение существует. Например, $x=e_{i}$.
Вектор $x=A^{-}a_i,$ по определению, является решением этой системы.

\pause

Из этого следует, что
\begin{eqnarray*}
Ax=AA^{-}a_i=a_i\ \ \forall i,
\end{eqnarray*}
что, в свою очередь, означает, что выполнено равенство $AA^{-}A=A$.
\bigskip
\pause

Пусть, теперь, система $Ax=y$ совместна и матрица $B$ удовлетворяет равенству
\begin{eqnarray*}
ABA=A.
\end{eqnarray*}
Тогда, умножая обе части этого равенства справа на $x$, получим
\begin{eqnarray*}
AB\underbrace{Ax}_y=\underbrace{Ax}_y\ \ \Rightarrow\ \ ABy=y,
\end{eqnarray*}
и, следовательно, вектор $x=By$ является решением системы $Ax=y$.
\bigskip
\pause

Это означает, по определению, что матрица $B$ является
обобщенно-обратной для $A$. То есть $B=A^{-}$.
 Лемма доказана. $$ \eqno \Box$$
 \end{frame}
 
\begin{frame}
Теперь рассмотрим вопрос существования обобщенно-обратной матрицы.
Справедлива следующая лемма.
\begin{lemma}{Лемма 6.}
Для всякой матрицы $A$ существует о.о. матрица $A^{-}$.
 \label{lemGIM2}
\end{lemma}
{\bf Доказательство леммы } Сперва рассмотрим случай,
когда $A$ - симметричная квадратная матрица. По теореме о
спектральном разложении, в этом случае, матрицу $A$ можно
представить в виде
\begin{eqnarray*}
A=P\Lambda P^T, \ \ PP^T=P^TP=I,
\end{eqnarray*}
$\Lambda$ - матрица, на диагонали которой стоят собственные числа
матрицы $A$,
 а внедиагональные элементы равны 0.
Без ограничения общности можно считать, что матрица $\Lambda$ имеет
вид:
\begin{eqnarray*}
\Lambda=\left(
          \begin{array}{cc}
            \Lambda_1 & 0 \\
            0 & 0 \\
          \end{array}
        \right),
\end{eqnarray*}
где $\Lambda_1$ имеет ненулевые диагональные элементы. Положим,
\begin{eqnarray*}
A^{-}=P\Lambda^{-}P^T,\ \  \Lambda^{-}=\left(
          \begin{array}{cc}
            \Lambda_1^{-1} & 0 \\
            0 & 0 \\
          \end{array}
        \right).
\end{eqnarray*}
\end{frame}

\begin{frame}
Проверим, что $AA^{-}A=A$:
\begin{eqnarray*}
AA^{-}A=P\Lambda P^TP\Lambda^{-}P^TP\Lambda P^T=P\Lambda \Lambda^{-}
\Lambda P^T=P\Lambda P^T=A.
\end{eqnarray*}
Следовательно, по лемме \ref{lemGIM1}, матрица $A^{-}$ является
обобщенно-обратной.

Для произвольной матрицы $A$ существует представление $A=B\Lambda
C$, где $B$ и $C$ - невырожденные матрицы, $\Lambda$ - диагональная.
Так же, как и в предыдущем случае, положим:
\begin{eqnarray*}
\Lambda=\left(
          \begin{array}{cc}
            \Lambda_1 & 0 \\
            0 & 0 \\
          \end{array}
        \right), \ \ \Lambda^{-}=\left(
          \begin{array}{cc}
            \Lambda_1^{-1} & 0 \\
            0 & 0 \\
          \end{array}
        \right),\ \ A^{-}=C^{-1}\Lambda^{-}B^{-1}.
\end{eqnarray*}
Несложно проверить, что для $A^{-}$ выполнено равенство $AA^{-}A=A$:
\begin{eqnarray*}
AA^{-}A=B\Lambda CC^{-1}\Lambda^{-}B^{-1}B\Lambda C=B\Lambda
\Lambda^{-} \Lambda C=B\Lambda C=A.
\end{eqnarray*}
Следовательно, по лемме \ref{lemGIM1}, матрица $A^{-}$ является
обобщенно-обратной.
 Лемма доказана. $$ \eqno \Box$$
 \end{frame}
\begin{frame}
Справедлива следующая теорема.

\begin{theorem}
Пусть $A^{-}$ какая-либо обобщенно-обратная для A матрица. И, пусть,
$H=A^{-}A$. Тогда,
\begin{itemize}
    \item[(a)] $H^2=H$, т.е. матрица $H$ идемпотентная;
    \item[(b)] $AH=A$ и ранг $A$ равен рангу $H$ и $tr(H);$
    \item[(c)] Общее решение системы уравнений $Ax=0$ имеет вид
    $x=(H-I)z$, где $z$-произвольный вектор;
    \item[(d)] Общее решение совместной системы уравнений $Ax=y$ имеет вид
    $x=A^{-}y+(H-I)z,$ где $z$ - произвольный вектор;
    \item[(e)] Произведение $Tx$ принимает единственное значение для
    всех $x$, удовлетворяющих системе $Ax=y$, тогда и только тогда,
    когда $TH=T.$
\end{itemize}
\label{TheoGIM1}
\end{theorem}
\end{frame}
\begin{frame}


{\bf Доказательство теоремы.} Пункты $(a)-(e)$ следуют
из леммы \ref{lemGIM1}. Прокомментируем вторую часть пункта $(b):$
Пусть ранг $A$ равен $k$.
\begin{eqnarray*}
tr A^{-}A=tr C^{-1}\Lambda^{-}B^{-1}B\Lambda C=
\\tr C^{-1}\Lambda^{-}\Lambda C =tr \Lambda^{-}\Lambda CC^{-1}=tr \Lambda^{-}\Lambda=tr
I_k=k,
\end{eqnarray*}
где $I_{k}$ диагональная матрица с $k$ единицами на главной
диагонали $k$- ранг матрицы $A$. Остальные ее элементы - нули. Мы
воспользовались известным результатом $tr AB=tr BA$ (для любых
согласованных матриц $A$, $B$).
 Теорема доказана. $$ \eqno \Box$$
\end{frame}
\begin{frame}
В случае, когда матрица $X$ имеет ранг меньший m, невозможно
несмещенно оценить все параметры регрессионной модели
$(Y,X\theta,\Sigma)$. Однако, можно оценить некоторые
параметрические функции  $\tau=T\theta$. Вид оценки и условия
оцениваемости определяются, в этом случае, в следующей теореме,
являющейся аналогом теоремы Гаусса-Маркова для вырожденной матрицы
$X^TX$.
\end{frame}
\begin{frame}
\begin{theorem}
Для классической регрессионной модели  и вектора
ошибок, удовлетворяющего условию $E\epsilon =0$,
\begin{itemize}
    \item[(1)] параметрическая функция $\tau=T\theta$ ($T\in R^{k\times m}$, $k \in \{1,\ldots,m\}$) оцениваема
тогда и только тогда, когда
\begin{eqnarray*}
T(X^TX)^{-}X^TX=T.
\end{eqnarray*}
   \item[(2)] Если выполнено условие $(1)$, оценка для $\tau$ имеет вид:
\begin{eqnarray*}
\widehat{\tau}= T(X^TX)^{-}X^TY,
\end{eqnarray*}
определена однозначно, и являются наилучшей линейной несмещенной
оценкой. Ковариационная матрица оценки $\widehat{\tau}$  имеет вид:
\begin{eqnarray*}
D_{\widehat{\tau}}=\sigma^2D,\ \ D=T(X^TX)^{-}T^T.
\end{eqnarray*}

\end{itemize}
\label{TheoGIM2}
\end{theorem}



\end{frame}
\begin{frame}{Нелинейные модели. Асимптотические свойства оценок параметров}
Вернемся снова к общему уравнению регрессии:
 \begin{align*}
	y_{i} = \eta(x_i,\theta) + \varepsilon_{i}, \qquad &i = 1, \dots, N.
 \end{align*}
 Пусть теперь $\eta(x,\theta)$ не обязательно линейная.
\pause 
 Будем предполагать, что

 Случайные ошибки $\varepsilon_1,\dots,\varepsilon_N$ независимы, одинаково распределены и обладают свойствами:   
 \begin{enumerate}
 \item $\mathbf{E}\varepsilon_i = 0$,\\
 \item $\mathbf{E}\varepsilon_i^2 = \sigma^2$. 
 
 \end{enumerate}
 \pause
Оценка вектора $\theta$ по методу наименьших квадратов определяется как решение экстремальной задачи:
\begin{gather*}
    \widehat\theta = \argmin_{\theta \in \mathbb{R}^m} \sum_{i = 1}^N (\eta(x_i, \theta) - y_i)^2.
\end{gather*}
\pause
Обсудим, какими асимптотическими свойствами обладает эта оценка.
\end{frame}

\begin{frame}
Введем предположения:
\begin{enumerate}
    \item
    Функция регрессии $\eta(x, \theta)$ непрерывна по $x$ и $\theta$;
	\pause
    \item
    Существует $\xi$ такой что $ \mathcal L(\xi_N) \Rightarrow \mathcal L(\xi)$, т.е. 
    $$
        \int_\mathcal X g(x) d\xi_N(x) \rightarrow  \int_\mathcal X g(x) d \xi(x), \qquad \forall g \in \mathrm C(\mathcal X);
    $$

	\pause
    \item Для любых $\theta$, $\overline\theta \in \Theta$
    $$
        \int_\mathcal X \left(\eta(x, \theta) - \eta(x, \overline \theta) \right)^2  d\xi(x) = 0 \; \Leftrightarrow \; \theta = \overline \theta;
    $$

	\pause
    \item
    $\eta \in C_\theta^2(\mathcal X \times \Theta)$;

	\pause
    \item
    Истинное значение параметра $\theta_u$ является внутренней точкой $\Theta$;

	\pause
    \item
    Матрица
    $$
        M(\xi, \theta) = \int_\mathcal X f(x, \theta) f^{T}(x, \theta)\, d\xi(x),
    $$
    где $
    f(x, \theta) = {\left[\frac{\partial \eta(x, \theta)}{\partial \theta_1}, \dots, \frac{\partial \eta(x, \theta)}{\partial \theta_m}\right]}^{T}
    $
    не вырождена при $\theta = \theta_u$.
\end{enumerate}
\end{frame}

\begin{frame}
Обозначим
\begin{gather}
\label{GLSE}
    \widehat\theta_N = \argmin_{\theta \in \mathbb{R}^m} \sum_{i = 1}^N (\eta(x_i, \theta) - y_i)^2.
\end{gather}
\begin{mytheorem}
    Если выполнены предположения 1--3, то последовательность МНК-оценок сильно состоятельна,
    т.е. при $N \to \infty$
    \begin{gather*}
        \widehat \theta_N \to \theta_u
    \end{gather*}
    с вероятностью 1, где $\widehat \theta_N$ определено формулой \eqref{GLSE}.

    Если дополнительно выполняются предположения 4--6, то последовательность случайных векторов 
    $$\sqrt N (\widehat \theta_N - \theta_u) \Rightarrow \mathcal{N}(\mathrm{0}, \Sigma),$$
     ковариационная матрица $\Sigma = \sigma^2 M^{-1}(\xi, \theta_u)$.
\end{mytheorem}
\end{frame}
\begin{frame}
\bigskip
Доказательство теоремы можно найти в работе

\bigskip

Jennrich, R. J. (1969). Asymptotic properties of non-linear least squares
estimators // Ann. Math. Stat., 40, 633–643.
\pause
\begin{enumerate}
\bigskip

    \item
Нелинейные модели обычно возникают как решения систем уравнений, описывающих поведение некоторого реального процесса. 
\bigskip
\pause
 \item
Параметры нелинейных моделей имеют, как правило, четкую предметную интерпретацию и находятся из анализа результатов эксперимнта.
\pause
\bigskip
 \item


Рассмотрим два примера.

 \end{enumerate}
\end{frame}

\end{document}
\begin{frame}{Пример 1, экспоненциальная модель}
Рассмотрим простейшее линейное дифференциальное уравнение с задачей Коши
\begin{align*}
\frac{d \eta(x)}{d x} = a \eta(x), \; \eta(0) = \eta_0.
\end{align*}

Такое уравнение при $a > 0$ может быть использовано для описания популяционной динамики во времени при наличии бесконечного количества ресурсов и отсутствии конкуренции.

Его решение есть простейшая экспоненциальная модель
\begin{align*}
\eta(x) = \eta_0 e^{a x}.
\end{align*}

\end{frame}

\begin{frame}{Пример 2, модель Моно}
Теперь рассмотрим более сложное уравнение
\begin{align*}
\frac{d \eta(x)}{d x} = \mu(x) \eta(x), \; \mu(x) = \theta_1\frac{s(x)}{s(x)+\theta_2},
\end{align*}
где
\begin{align*}
s(x) - s_0 = \frac{\eta_0 - \eta(x)}{\theta_3}, \; \eta(0) = \eta_0, \; s(0) = s_0. 
\end{align*}
\end{frame}
\begin{frame}
Решением этого уравнения на интервале $x \in [0,T]$ является модель Моно, которая задается неявно следующими уравнениями:
  \begin{align}
	x = &\frac{1}{\theta_{1}} \left[ (1+b) \log{\frac{\eta(x)}{\eta_0}} - b \log{\frac{a - \eta(x)}{a - \eta_0}} \right];\\
	 &a = s_0 \theta_3 +\eta_0,\; b = \frac{\theta_{2} \theta_{3}}{a}.
  \end{align}
\pause

\bigskip

Изучению этой модели посвящена работа

\bigskip

H. Dette, V. Melas, A. Pepelyshev, N. Strigul (2003). Efficient design of experiments in the Monod model. Journal of the Royal Statistical Society: Series B (Statistical Methodology, Volume 65, Issue 3

\end{frame}

\end{document}